\documentclass[10pt,twoside]{uz_kgu}
\usepackage[T2A]{fontenc}
\usepackage[utf8]{inputenc}
\usepackage[russian,english]{babel}
%\inputencoding{cp866}%(DOS)
%\inputencoding{koi8-r}%(UNix)
%\inputencoding{cp1251}%(Win)
\begin{document}
\renewcommand{\refname}{\small {Литература}} \renewcommand{\abstractname}{Аннотация} \renewcommand{\figurename}{Рис.}
\renewcommand{\proofname}{ {\hskip\parindent \bf Доказательство. }}
\UDK{???????}
\ArticleNAME{ПРИМЕР ОФОРМЛЕНИЯ СТАТЬИ  \\ НАЗВАНИЕ СТАТЬИ ЗАГЛАВНЫМИ БУКВАМИ}
\ArticleAUTHOR{А.А.~Автор1, В.В.~Автор2}
\ArticleHEAD{Название статьи для колонтитула}
\ArticleAUTHORHEAD{А.А.~Автор1, В.В.~Автор2}
\makeabstitle
\begin{abstract}
Текст текст текст текст текст текст текст текст текст текст текст текст текст текст текст текст текст текст текст текст текст текст текст текст текст текст текст текст текст текст текст текст текст текст текст текст текст текст текст текст текст текст текст текст текст текст текст текст текст текст текст текст.
\vspace{3pt}
\textbf{Ключевые слова:} (до 10 словосочетаний)
\end{abstract}
\vspace{\baselineskip}\hrule
\section{Пункт1}
Текст статьи должен быть подготовлен средствами издательской системы \LaTeXe \ с использованием стиля  uz\_kgu.cls. В случае использования нестандартных стилевых файлов автор обязан предоставить необходимые стилевые файлы. Изменения стандартных стилевых файлов недопустимы.
Рекомендуется вставлять текст статьи непосредственно <<поверх>> текста этого шаблона.
В статье нумеруются лишь те формулы, на которые по тексту есть ссылки.
Список литературы составляется по порядку ссылок в тексте (где номер ссылки указывается в квадратных скобках) и приводится в конце статьи. Ссылки на неопубликованные работы не допускаются. Пример оформления ссылки см. ниже.
\section{Пункт2}

Основной текст статьи.
Рисунки могут быть подготовлены в любом графическом редакторе и предоставлены в формате EPS. Рисунки должны быть черно-белые или полутоновые (серой градации).
 Пример вставки файла с рисунком в формате EPS.
%
\begin{figure}[!tbh]
	\centering
	\includegraphics[clip,width=50mm]{fig.eps}
\caption{Пример вставки рисунка в формате EPS}\label{tempvoid1}
\end{figure}
Для удобства верстки, необходимо, чтобы каждый автор все свои метки ({teo1}, {eqv1}, {tab}, {bib1} и т.~д.(см. ниже)) называл как-либо "уникально". Например: {Ivanovbib1}, {Ivanovtheo1} и т.~п.
Сноски нужно оформлять следующим образом \footnote{Пример оформления сноски}.
\section{Пункт3}
\begin{Theorem} \label{tempteo1} Формулировка теоремы.
\end{Theorem}
\begin{proof}
Текст доказательства. Пример использования формулы с номером:
\begin{equation} \label{tempeqv1}
		a(x)=b^3.
\end{equation}
\begin{equation} \label{tempeqv2}
		c(x)=d^2.
\end{equation}
\begin{equation} \label{tempeqv3}
		f(x)=x^a.
\end{equation}
..........
Теорема доказана.
\end{proof}
Возможны ссылки на формулы. Например, здесь мы сослались на формулу (\ref{tempeqv1}) из предыдущей теоремы \ref{tempteo1}. Пример ссылок на несколько формул (\ref{tempeqv1}), (\ref{tempeqv2}) или же (\ref{tempeqv1})--(\ref{tempeqv3}).
Пример использования формулы без номера:
	$$F(x)=x-4.$$
Пример задания ссылок на используемую литературу~\cite{tempbib1, tempbib3} или~\cite{tempbib2}.
\vspace{\baselineskip} Благодарности ...

%Может отсутствовать
\vspace{\baselineskip} Работа выполнена при финансовой поддержке ... %Может отсутствовать
\vspace{\baselineskip}
\renewcommand{\abstractname}{Summary}
\begin{abstract}

{\it A.A.~Author1, В.В.~Author2 (Фамилии авторов на английском языке).} Название статьи на английском языке.
\vspace{3pt}
 Text text text text text text text text text text text text text text
 text  text text text text text text text text text text text text text text text text text text text text text text text text text 
text text text text text text text text text text text text

 text.
 \vspace{3pt}
\textbf{Key words:} (ключевые слова на англ. яз.)
\end{abstract}
%	Пример списка литературы
\References{
\bibitem{tempbib1}
{\it Баренблатт Г.И., Ентов В.М., Рыжик В.М.\/} Движение жидкостей и газов в природных пластах. -- М.: Недра, 1984. -- 211~с. \bibitem{tempbib2}
{\it Мироненко В.А., Петров Н.С. \/} Загрязнение подземных вод углеводородами // Геоэкология. -- 1995. -- №~1. -- С.~3--27. \bibitem{tempbib3}
{\it Pullan  A.J.\/} The  quasilinear   approximation   for unsaturated porous  media  flow // Water  Resourses Research. -1990. -- V.~26, No~3. -- P.~1219--1234.
\bibitem{tempbib4}
{\it Баренблатт Г.И., Ентов В.М., Рыжик В.М.\/} Движение жидкостей и газов в природных пластах. -- М.: Недра, 1984. -- 211~с. \bibitem{tempbib5}
{\it Мироненко В.А., Петров Н.С. \/} Загрязнение подземных вод углеводородами // Геоэкология. -- 1995. -- №~1. -- С.~3--27. }
{\small \vspace{\baselineskip}\hrule \vspace{3pt}
\par
Сведения о каждом из авторов статьи
\par
{\bf Фамилия, имя и отчество} -- ученая степень, ученое звание, должность и организация, в которой работает автор
\par
E-mail: {\it электронный адрес} }
\end{document}
